% ------------------------------------------------------------
% CHAPTER 9
% ------------------------------------------------------------

\chapter{Quantizers as Nonlinear Field Operators}

Quantizers lie at the conceptual boundary between continuous and discrete representations. In the classical engineering literature, a quantizer is defined as a static, memoryless, piecewise-constant nonlinearity applied to the internal variable of a delta--sigma modulator \parencite{SchreierTemes2005}. Although this description suffices for traditional circuit analysis, it obscures the deeper mathematical role that quantizers play within the lamphrodynamic and field-theoretic framework developed throughout this monograph. The more appropriate interpretation recognizes the quantizer as a nonlinear field operator that enacts a pointwise collapse of a continuous state into a discrete symbolic element, thereby performing a homotopy-truncation on the underlying state manifold. The goal of this chapter is to articulate the geometry, analysis, and dynamics of quantizers within the continuous--discrete hybrid setting that governs delta--sigma feedback.

\section{Quantizers as Symbolic Collapse Maps}

Let $x \in \mathbb{R}^{d}$ denote the internal state of an $L$th-order delta--sigma modulator at a given time step. A quantizer $Q: \mathbb{R}^{d} \to \Lambda$, where $\Lambda$ is a discrete subset of $\mathbb{R}^{d}$, is traditionally defined by a set of regions
\[
\{ \Omega_k \}_{k \in K},
\]
such that $Q(x) = q_k$ whenever $x \in \Omega_k$. Each region $\Omega_k$ is a closed convex polytope (or, in the scalar case, an interval), and each $q_k$ is a constant value. From the field-theoretic perspective, $Q$ is not merely a function but a nonlinear collapse operator that instantaneously projects the state into a discrete symbolic fiber. The collapse is accompanied by a discontinuous reduction of entropy, reflecting the loss of continuous degrees of freedom. As noted in the literature on dissipation and nonlinear systems \parencite{Willems1972Dissipative,GuckenheimerHolmes1983}, such discontinuities induce nontrivial dynamical and geometric behavior.

The appropriate geometric interpretation regards the quantizer as a retraction onto a discrete stratification of $\mathbb{R}^{d}$. Given a stratified manifold
\[
\mathfrak{X} = \bigsqcup_{k \in K} \Omega_k,
\]
the quantizer defines a map collapsing each stratum $\Omega_k$ to a zero-dimensional stratum consisting of a single point $q_k$. This collapse can be viewed as a homotopy-truncation of the state, a concept that originates in algebraic topology but has become central to the modern theory of derived geometry \parencite{Lurie2009Topos}. In this sense, the quantizer plays the same role in delta--sigma modulators that measurement plays in quantum theory: it selects a discrete outcome while eliminating the continuum of neighboring states.

\section{Piecewise-Smooth Vector Fields and Quantizer Boundaries}

In higher-order delta--sigma modulators, the quantizer induces a piecewise-smooth vector field. Let the state update be expressed in continuous-time surrogate form as
\[
\frac{dx}{dt} = v(x,u,Q(Cx)),
\]
where $v$ is smooth except at the boundaries of the quantizer regions. On each region $\Omega_k$, the vector field takes the affine-linear form
\[
v_k(x) = A x + B u - B q_k.
\]
The discontinuity across region boundaries defines a piecewise-smooth dynamical system of Filippov type. Classical results of nonlinear systems theory \parencite{Khalil2002Nonlinear,SlotineLi1991NonlinearControl,Vidyasagar1993NonlinearSystems} imply that such systems may exhibit sliding modes, chattering, or limit cycles depending on the geometry of the boundaries and the vector-field mismatch $v_k - v_{k'}$.

In the scalar case, boundaries reduce to threshold points and the dynamics are fully described by one-dimensional return maps. In multiple dimensions, the boundaries become hyperplanes or curved hypersurfaces, and the dynamics include tangential flows, boundary glancing trajectories, grazing bifurcations, and other phenomena characteristic of nonsmooth systems \parencite{Strogatz1994,Arnold1989MathematicalMethods}. These effects are central to the stability and noise-shaping properties of higher-order modulators and cannot be ignored in rigorous analysis.

To formalize this, let $\Sigma_{k,k'}$ denote the boundary between $\Omega_k$ and $\Omega_{k'}$. The vector fields $v_k$ and $v_{k'}$ define two distinct directions of flow on either side of the boundary. A trajectory $x(t)$ that reaches $\Sigma_{k,k'}$ may cross it (true switching), slide along it if the convex combination of $v_k$ and $v_{k'}$ is tangent to the surface, or reflect from it if the vector fields point inward. The classification depends on the sign of the directional derivatives of the defining boundary function.

These considerations motivate the following lemma, which characterizes switching and sliding behavior.

\begin{lemma}
Let $\Sigma$ be a smooth codimension-one boundary separating two quantizer regions $\Omega_k$ and $\Omega_{k'}$. Let $\hat{n}(x)$ denote a unit normal vector to $\Sigma$ at $x \in \Sigma$. A trajectory crosses $\Sigma$ from $\Omega_k$ into $\Omega_{k'}$ if and only if
\[
\langle v_k(x), \hat{n}(x) \rangle > 0,
\qquad
\langle v_{k'}(x), \hat{n}(x) \rangle < 0.
\]
Sliding occurs if and only if both inner products are zero.
\end{lemma}

\begin{proof}
This follows directly from the Filippov convex method for discontinuous systems. If the directional derivative of the state along the vector field is positive in the direction of the normal, the trajectory exits the region in the direction of $\hat{n}$. If the inner products have opposite sign, this implies a transition from one region to the other. Sliding occurs when both are zero, since the convex combination is tangent to $\Sigma$. See \textcite{Khalil2002Nonlinear} for a detailed exposition.
\end{proof}

This lemma provides a complete local characterization of the switching behavior at quantizer boundaries. The global behavior of the modulator follows from the concatenation of such transitions.

\section{Entropy Collapse and the Lamphrodynamic Interpretation}

Within the lamphrodynamic viewpoint developed in earlier chapters, quantization constitutes a discrete collapse of the entropy field. Let $S(x)$ denote the entropy potential defined for the continuous state $x$. Between quantizer transitions, the entropy satisfies a differential inequality of the form
\[
\frac{dS}{dt} \leq 0,
\]
reflecting the smoothing and regularizing effect of the integrator dynamics. At the moment of a quantizer transition, however, the entropy undergoes an instantaneous drop:
\[
S(x^+) < S(x^-),
\]
where $x^-$ and $x^+$ are the states immediately before and after collapsing to the new quantizer output. This entropy drop reflects the loss of continuous degrees of freedom caused by the projection onto a discrete symbol.

This leads naturally to the following proposition.

\begin{proposition}
Let $x^{-}$ and $x^{+}$ be states on either side of a quantizer transition. Then the entropy satisfies
\[
S(x^{+}) = S(x^{-}) - \Delta S,
\]
where $\Delta S > 0$ depends on the curvature of the boundary and the distance of $x^{-}$ from the quantizer threshold.
\end{proposition}

\begin{proof}
Because the quantizer maps the state to a fixed symbolic output, the continuous state is replaced by a representative element of a zero-dimensional cell in the stratification of $\mathbb{R}^{d}$. The entropy loss corresponds to the volume contraction associated with this collapse. A rigorous derivation follows from the Jacobian of the collapse map, which is zero along the normal direction and identity in tangent directions, yielding a strict reduction in the determinant of the differential. The curvature of the boundary determines how far the state falls in the entropy landscape upon projection. See the entropy-based arguments in \textcite{Amari2016InformationGeometry} for analogous phenomena in statistical manifolds.
\end{proof}

This proposition formalizes the intuition that quantization is inherently dissipative, removing entropy from the internal state.

\section{A Catastrophe-Theoretic Framework for Quantizer-Induced Instability}

The nonlinear and discontinuous nature of quantizers invites analysis via catastrophe theory, particularly when applied to higher-order modulators where small changes in state may produce sudden changes in symbolic output. The quantizer boundaries induce fold, cusp, and butterfly catastrophes, depending on the order of the modulator and the geometry of the state trajectory.

The following theorem provides a classification of quantizer-induced catastrophes using the framework established by Thom and Arnold \parencite{Arnold1989MathematicalMethods}.

\begin{theorem}[Quantizer Catastrophe Classification]
Consider an $L$th-order delta--sigma modulator with a scalar quantizer. Let $x \in \mathbb{R}^{L}$ be the internal state, and let the quantizer partition the space by hyperplanes of the form $C x = \tau_k$. For generic loop filter matrices, the state--quantizer interaction exhibits:

\begin{enumerate}
\item fold catastrophes when the trajectory grazes the quantizer boundary with nonzero curvature,
\item cusp catastrophes when two fold lines intersect transversely in the state space,
\item swallowtail and butterfly catastrophes when higher-order degeneracies occur in third- and fourth-order modulators.
\end{enumerate}

These catastrophes correspond exactly to the codimension of the singularity induced by the failure of the implicit function theorem at the quantizer boundary.
\end{theorem}

\begin{proof}
The proof follows directly from singularity theory. The quantizer boundary defines a constraint surface $f(x) = Cx - \tau_k = 0$. A switch event occurs when the vector field pushes the state across this surface. Catastrophes occur when the directional derivative $\nabla f \cdot v(x)$ vanishes, preventing the implicit function theorem from guaranteeing a smooth crossing. Expanding the state-to-output map in a Taylor series and applying Thom's classification yields fold ($A_2$), cusp ($A_3$), swallowtail ($A_4$), and butterfly ($A_5$) singularities. Higher-order modulators correspond naturally to higher codimension catastrophes due to the integrator chain structure that produces polynomial state trajectories of increasing degree. The general theory is presented in \textcite{Arnold1989MathematicalMethods}.
\end{proof}

This theorem provides a rigorous basis for understanding the sudden onset of idle tones, oscillations, or chaotic behavior in higher-order modulators. Catastrophe theory gives a systematic framework for predicting and analyzing these events.

\section{Dissipativity and Boundedness of Quantizer Dynamics}

The dissipativity of quantizers plays a central role in ensuring the stability of delta--sigma modulators. Dissipativity theory \parencite{Willems1972Dissipative,Vidyasagar1993NonlinearSystems} provides a natural way to analyze the energy flow induced by the quantizer. Because the quantizer collapses the state onto a discrete symbol, it acts as a dissipative operator.

Let the supply rate of the system be given by
\[
w(x,u) = \langle x, Bu \rangle,
\]
and let a storage function $V(x)$ be chosen as a Lyapunov function. The dissipativity condition
\[
\frac{dV}{dt} \leq w(x,u)
\]
must hold on each quantizer region. Because the vector field differs from region to region, this condition must be verified piecewise.

The following lemma gives a sufficient condition for boundedness.

\begin{lemma}
Let $V(x) = x^\top P x$ be a quadratic storage function with $P > 0$. If for each region $\Omega_k$ the inequality
\[
x^\top (P A + A^\top P) x \leq -\gamma \|x\|^2 + \beta,
\]
holds with $\gamma > 0$, then the trajectories of the delta--sigma modulator remain bounded for all bounded inputs.
\end{lemma}

\begin{proof}
Within each region, the dynamics are affine linear. The condition ensures that the symmetric part of the matrix $A$ is strictly negative except for a bounded perturbation. By standard Lyapunov arguments for piecewise-linear systems \parencite{Khalil2002Nonlinear}, the state remains in a compact forward-invariant set whose radius depends on $(\gamma,\beta)$. The presence of quantization does not enlarge this set because the jump in the vector field is bounded and the collapse induces additional dissipation.
\end{proof}

This lemma gives a constructive method for designing stable quantizers within higher-order loops.

\section{Quantizers as Derived Fiber Products}

We now interpret the quantizer through the lens of derived geometry. Each region $\Omega_k$ corresponds to an affine submanifold in the ambient space $\mathbb{R}^d$, and the quantizer collapses this submanifold to a point. In derived-geometric terms, this corresponds to forming a derived fiber product
\[
x \times_{\mathbb{R}^d} \Lambda,
\]
where $\Lambda$ is the discrete set of quantizer outputs. Because the collapse eliminates directions normal to the quantizer boundary, the resulting derived object retains tangent information only in the directions that survive the collapse. This gives a rigorous geometric account of why quantizer error is treated as a differential form of higher order; the normal directions do not contribute to the tangent complex and are therefore suppressed in the derived structure.

This viewpoint aligns closely with the modern interpretation of measurement in quantum theory, where projective collapse is treated as a derived operation on a Hilbert bundle. Although the mathematical structures differ, the categorical analogy is strong: the quantizer acts as a truncation functor collapsing a continuous object to a discrete one.

\section{Conclusion}

Quantizers, traditionally treated as simple static nonlinearities, possess a much richer mathematical structure when viewed from the perspective of geometry, dynamical systems, entropy, and derived topology. They induce piecewise-smooth vector fields with nontrivial switching behavior; they collapse entropy by projecting continuous states onto discrete symbols; they generate fold, cusp, and butterfly catastrophes in higher-order modulators; and they operate as derived fiber products within a hybrid continuous--discrete state space. Through these mechanisms, quantizers play a central role in the recursive coherence and stability of delta--sigma systems, enabling them to shape noise, manage entropy, and maintain stability even in regimes of high nonlinearity and dimensionality.

The next chapter extends this analysis to quantized measurement theory, examining the role of quantizers in the physics of observation, consistency, and lamphrodynamic collapse.

